
(G) \uline{Il
existe un} @@M0@@-\uline{schéma lisse} @@M1@@, \uline{factorisant} @@M2@@, \uline{tel que pour toute immersion fermée} : @@M3@@ \uline{on puisse trouver un} @@M4@@-\uline{schéma lisse} @@M5@@, \uline{avec} @@M6@@ (\uline{cette condition est en particulier remplie si} @@M7@@ \uline{est quasi-projectif et si} @@M8@@ \uline{a un faisceau ample}, \uline{en prenant pour} @@M9@@ \uline{un} @@M10@@).

\uline{Alors} @@M11@@ \uline{est d'intersection complète}, \uline{lissifiable}, \uline{et on a un triangle distingué} :
@@M12@@

Soit @@M13@@ un @@M14@@-schéma lisse factorisant @@M15@@. On considère le diagramme commutatif :
@@M16@@

Alors l'immersion @@M17@@ est régulière, car déduite de @@M18@@ par un changement de base plat. Par suite @@M19@@ est une immersion régulière factorisant @@M20@@, qui est donc lissifiable et d'intersection complète. De plus, on a la suite exacte des faisceaux conormaux :
@@M21@@
en notant @@M22@@ les faisceaux définissant respectivement les immersions @@M23@@. Or, @@M24@@ étant plat, @@M25@@. La suite peut donc s'écrire :
@@M26@@

Par ailleurs, @@M27@@ et @@M28@@ étant lisses sur @@M29@@, on a la suite exacte :
@@M30@@
